% GENERATED BY final_submission/scripts/make_paper_tables.py
% Do not edit by hand: regenerate with
%   python final_submission/scripts/make_paper_tables.py \
%       --data-dir final_submission/data \
%       --out final_submission/paper/tables.tex

\begin{table*}[t]
\centering
\scriptsize
\caption{Measured video-flow latency ($h1\!\to\!h4$) in every phase of the experiment (S6, with the lab stopped, yields no samples and is omitted), and the protected voice flow ($h3\!\to\!h4$) measured in the same snapshots. Each phase is five independent snapshots of ten ICMP samples each; \emph{Snaps.} is the number of snapshots, not the per-snapshot sample count ($10$). \emph{Mean} is the mean of the per-snapshot means; \emph{Peak} is the largest single ping sample seen in the phase, not the largest snapshot mean. The agent's own before/after measurements (Tab.~\ref{tab:outcomes}) are taken independently of the phase sampling and can differ from it. From \texttt{final\_submission/data/measurements.csv}.}
\label{tab:phases}
\begin{tabular}{llrrrrrr}
\toprule
Phase & Intent & Snaps. & Mean & Median & Std.~dev. & Peak & Voice \\
\midrule
S1 healthy & 50~ms & 5 & 17.8 & 17.2 & 1.2 & 46.0 & 0.6 \\
S2 degraded & 50~ms & 25 & 98.1 & 97.3 & 2.2 & 192.5 & 0.4 \\
S3 degraded & 22~ms & 25 & 97.4 & 97.3 & 0.5 & 122.2 & 0.4 \\
S4 healthy & 15~ms & 5 & 17.6 & 17.2 & 0.7 & 34.8 & 0.5 \\
S5 healthy & -- & 5 & 17.3 & 17.3 & 0.1 & 19.7 & 0.4 \\
S7 degraded+load & 50~ms & 25 & 97.8 & 97.7 & 0.5 & 111.3 & 0.4 \\
S8 degraded, no flag & 50~ms & 25 & 97.3 & 97.3 & 0.1 & 98.7 & 0.4 \\
S1' restart & -- & 5 & 17.3 & 17.3 & 0.1 & 18.0 & 0.4 \\
\bottomrule
\end{tabular}
\end{table*}

\begin{table*}[t]
\centering
\scriptsize
\caption{Outcome of every agent run. ``Before'' and ``after'' are the latencies the verification gate measured; ``--'' means the run never changed the network. The gate column is the number of the nine deterministic safety checks that passed. From \texttt{final\_submission/data/scenario\_summary.csv}.}
\label{tab:outcomes}
\begin{tabular}{llrrrl}
\toprule
Scenario & Outcome & Before & After & Installed path & Gates \\
\midrule
S1 healthy & \texttt{intent\_satisfied} & -- & -- & \texttt{--} & -- \\
S2 degraded & \texttt{intent\_satisfied (5 trials)} & 99.3 [97.2--107.4] & 26.7 [25.3--31.9] & \texttt{s1-s3-s4} & 9/9 \\
S3 degraded & \texttt{partial\_improvement (5 trials)} & 97.3 [97.2--97.4] & 25.4 [25.3--25.4] & \texttt{s1-s3-s4} & 9/9 \\
S4 healthy & \texttt{no\_better\_path} & -- & -- & \texttt{--} & -- \\
S5 guardrail & \texttt{guardrail\_blocked} & -- & -- & \texttt{--} & -- \\
S6 lab offline & \texttt{lab\_offline} & -- & -- & \texttt{--} & -- \\
S7 load & \texttt{intent\_satisfied (5 trials)} & 97.7 [97.0--98.3] & 25.2 [25.0--25.3] & \texttt{s1-s3-s4} & 9/9 \\
S8 degraded, no flag & \texttt{intent\_satisfied (5 trials)} & 97.3 [97.3--97.3] & 25.3 [25.3--25.5] & \texttt{s1-s3-s4} & 9/9 \\
\bottomrule
\end{tabular}
\end{table*}

\begin{table*}[t]
\scriptsize
\centering
\caption{Outcome distribution per scenario. The active-rerouting scenarios (S2 recovery, S3 partial improvement, S7 offered load, S8 no-oracle) are repeated as five independent trials, each from a freshly built lab; the other scenarios are run once. The latency column is the post-action latency the verification gate measured; the wall-clock column is the end-to-end closed-loop duration (pre-action snapshot to verified outcome), min--max over the trials. From \texttt{final\_submission/data/scenario\_aggregates.csv}.}
\label{tab:aggregates}
\begin{tabular}{llrrll}
\toprule
Scenario & Trials & Outcome(s) & After (min--max) & Wall-clock (min--max) & Path \\
\midrule
S1 healthy & 1 & \texttt{intent\_satisfied (1)} & -- & 1.9--1.9 s & \texttt{--} \\
S2 degraded & 5 & \texttt{intent\_satisfied (5)} & 25.3--31.9 ms & 5.7--6.2 s & \texttt{s1-s3-s4} \\
S3 degraded & 5 & \texttt{partial\_improvement (5)} & 25.3--25.4 ms & 5.8--6.2 s & \texttt{s1-s3-s4} \\
S8 degraded, no flag & 5 & \texttt{intent\_satisfied (5)} & 25.3--25.5 ms & 5.8--5.8 s & \texttt{s1-s3-s4} \\
S4 healthy & 1 & \texttt{no\_better\_path (1)} & -- & 3.8--3.8 s & \texttt{--} \\
S5 guardrail & 1 & \texttt{guardrail\_blocked (1)} & -- & 0.0--0.0 s & \texttt{--} \\
S6 lab offline & 1 & \texttt{lab\_offline (1)} & -- & 0.1--0.1 s & \texttt{--} \\
S7 load & 5 & \texttt{intent\_satisfied (5)} & 25.0--25.3 ms & 6.2--6.2 s & \texttt{s1-s3-s4} \\
\bottomrule
\end{tabular}
\end{table*}

\begin{table}[t]
\centering
\caption{Path selection under the degraded bottleneck. Both rows rank the \emph{same} candidate set the lab produced, so the comparison isolates the ranking rule. The measured latency is what the live network recorded on the selected path. From \texttt{final\_submission/data/ablations.csv}.}
\label{tab:ranking}
\begin{tabular}{llr}
\toprule
Ranking rule & Path & Latency \\
\midrule
system (congestion-aware) & \texttt{s1-s3-s4} & 25.4~ms \\
ablation (shortest path) & \texttt{s1-s2-s4} & 97.2~ms \\
\bottomrule
\end{tabular}
\end{table}

\begin{table}[t]
\centering
\caption{Recorded measurements on which the two violation rules disagree, one row per direction. The system treats a flow as violated when at least 25\% of its individual ping samples exceed the target; the mean-only ablation compares only the mean. \emph{System / mean-only} reports each rule's verdict on the intent (yes = violated). \texttt{yes / no} is the oscillation case the sample-distribution rule is built to catch; \texttt{no / yes} is the outlier-tolerance case (one high sample, below $\gamma$). Rows are the closest calls in the recorded data (smallest $|\bar{s}-\tau|$), from a millisecond sweep of $\tau$ from $5$ to $200$~ms, not an estimate of the deployed operating point. From \texttt{final\_submission/data/ablations.csv}.}
\label{tab:violation}
\tiny
\begin{tabular}{@{}l r r l >{\raggedright\arraybackslash}p{0.34\columnwidth}@{}}
\toprule
Phase & Target & Mean & Sys / mean & Samples (ms) \\
\midrule
S2 degraded & 98 & 98.1 & no / yes & {\ttfamily\scriptsize 105.0, 97.4, 97.2, 97.2, 97.4, 97.3, 97.3, 97.2, 97.3, 97.4} \\
S7 degraded+load & 98 & 98.0 & yes / no & {\ttfamily\scriptsize 96.8, 98.6, 97.6, 99.1, 99.0, 97.6, 98.4, 98.0, 97.9, 96.9} \\
\bottomrule
\end{tabular}
\end{table}

\begin{table}[t]
\centering
\caption{Acting when there is nothing to act on. The system ends a healthy run with no flow rules installed; the always-act variant takes the same code path the agent uses when a target is unreachable, so its cost is measured rather than asserted. From \texttt{final\_submission/data/ablations.csv}.}
\label{tab:alwaysact}
\scriptsize
\begin{tabular}{llr}
\toprule
Policy & Outcome & Healthy flow \\
\midrule
system & \texttt{no-op (nothing installed)} & 17.6~ms, unchanged \\
always act & \texttt{rerouted} & 17.6~$\to$~25.7~ms (+46\%) \\
\bottomrule
\end{tabular}
\end{table}

